\chapter{ESPECIFICACIONES Y EJECUCIÓN}
\section{Equipos y materiales}
El conductor es AAAC 6201-T81 de 50 mm$^2$ nominales. Los postes se agrupan en 13, 15 y 17 m con capacidad transversal tipo de 600 daN para la verificación presentada. Las retenidas se especifican con resistencia mínima de rotura de 60 kN. Aisladores, crucetas, grapas y pernos corresponden a la función de paso, ángulo, retención o terminal indicada por apoyo.

Los transformadores se modelan con 250 kVA netos en sitio, considerando la altitud como condición de prestación. La maniobra y protección se organizan en los extremos de línea y transformación. La tensión de aislamiento y relación física de equipos dependen de la alternativa nominal; los cálculos no equiparan una relación de 22 kV con otra de 22.9 kV.
\section{Montaje y control de calidad}
La ejecución física comprende replanteo, liberación de accesos y servidumbres, cimentaciones, izado, instalación de armados, tendido por poleas, regulación de flechas, continuidad eléctrica, puesta a tierra y pruebas de maniobra/protección. La tabla de catenaria es una base de cálculo a temperaturas y cargas definidas; no sustituye la tabla de tendido asociada al material efectivamente suministrado.

Los controles comprenden certificados de conductor y herrajes, identificación de postes, integridad de aisladores, torque de conexiones, continuidad de tierras, relación y polaridad de TC/TT, aislamiento, relación de transformación y pruebas funcionales de disparo e interbloqueos. El sistema aislado exige comprobar la respuesta de frecuencia y tensión frente a escalones de carga y pérdida de generación.
\section{Referencias normativas y documentales}
Se distingue el marco de instalaciones de suministro de las instalaciones interiores. El CNE Suministro sustenta los criterios de línea y la interfaz de equipos; el CNE Utilización corresponde a instalaciones interiores y servicios auxiliares dentro de su ámbito. El RNE y las disposiciones de seguridad y servidumbres se consideran según la naturaleza de la obra.

Las referencias documentales utilizadas en el desarrollo son:
\begin{enumerate}
\item MINEM, Código Nacional de Electricidad -- Suministro 2011, RM 214-2011-MEM/DM. Copia pública consultada en la base de diseño.
\item MINEM, RD 018-2003-EM/DGE, Bases para el diseño de líneas y redes primarias para electrificación rural.
\item MINEM, RD 024-2003-EM/DGE, Especificaciones técnicas de soportes normalizados para líneas y redes primarias. Las láminas adjuntas son referencias documentales vectorizadas.
\item INEI, Directorio Nacional de Centros Poblados 2017, tomo 4, registros del ámbito de Huancasayani.
\item Datos de cartografía OSM, catastro aportado y DEM público empleados en el corredor de gabinete.
\end{enumerate}
La referencia de armados no acredita automáticamente una adaptación a poste CAC ni una lista de fabricación. Las hojas conservan su título y lista de componentes, y se separan de las láminas propias de la línea.
